\section{Related Work}
\label{sec:related}

Our work sits at the intersection of three research areas: goal drift in LLM agents, evolutionary prompt optimization, and co-evolutionary defense against prompt manipulation.

\para{Goal drift and task persistence.}
\citet{arike2025goaldrift} formally define goal drift in autonomous LLM agents, distinguishing between drift by commission (actively pursuing misaligned goals) and drift by omission (failing to divest from misaligned positions).
Their key finding---that ``strong elicitation'' prompts with emphatic goal statements significantly reduce drift---directly motivates our work: if prompt wording matters, can evolution find better wordings than humans?
\citet{abdelnabi2024tasktracker} detect task drift using activation deltas between hidden states before and after processing external data, achieving ROC AUC $\geq$0.994 across multiple models.
Their activation-based approach suggests an alternative fitness signal for future evolution work, though it requires white-box model access.
Unlike these papers, which measure drift but do not optimize against it, we use drift metrics as a fitness function for evolutionary prompt search.

\para{Evolutionary prompt optimization.}
\evoprompt \citep{guo2024evoprompt} applies genetic algorithms (GA) and Differential Evolution (DE) to prompt optimization using LLMs as crossover and mutation operators.
DE is more robust to poor initialization and generates more diverse prompts, making it our preferred variant since we do not know \emph{a priori} what good goal-persistent prompts look like.
\promptbreeder \citep{fernando2024promptbreeder} extends this with self-referential evolution, simultaneously evolving task-prompts and mutation-prompts.
\opro \citep{yang2023opro} takes an alternative approach, presenting LLMs with instruction-score pairs and asking for improvements without explicit evolutionary operators.
Additional work includes \textsc{GPS} \citep{xu2022gps} for genetic prompt search, \textsc{SPELL} \citep{gao2023spell} for semantic prompt mutations, \textsc{SPRIG} \citep{srinivasan2024sprig} for system prompt optimization, \textsc{TARE} \citep{tare2025} for robustness-aware evolution, and \textsc{GAAPO} \citep{gaapo2025} for GA-based prompt optimization.
All of these methods optimize for task accuracy; we are the first to apply evolutionary prompt optimization with a goal-persistence fitness function.

\para{Co-evolutionary defense.}
\aegis \citep{aegis2025} uses a GAN-inspired co-evolutionary framework where attacker and defender prompt agents alternate optimization.
Their evolved defense prompts achieve 84\% true positive rate versus 64\% for human-crafted prompts, demonstrating that evolution can discover robustness properties that humans miss.
\textsc{EvoRefuse} \citep{evorefuse2025} applies evolutionary optimization to mitigate LLM over-refusal, the opposite extreme of goal drift.
Building on these ideas, our work replaces the attacker/defender framing with a distraction/persistence framing, and demonstrates that while evolution discovers richer prompt structures, the performance advantage depends on task difficulty.
